\begin{ex}
Viết số 345 thành tổng: $300 + 40 + \dots$ Số còn thiếu là bao nhiêu?
\shortans{5}
\loigiai{
Số 345 gồm 300, 40 và 5.
}
\end{ex}

\begin{ex}
Viết số 567 thành tổng: $500 + \dots + 7$. Số còn thiếu là bao nhiêu?
\shortans{60}
\loigiai{
Số 567 gồm 500, 60 và 7.
}
\end{ex}

\begin{ex}
Viết số 728 thành tổng: $\dots + 20 + 8$. Số còn thiếu là bao nhiêu?
\shortans{700}
\loigiai{
Số 728 gồm 700, 20 và 8.
}
\end{ex}

\begin{ex}
Số 408 viết thành tổng các trăm và đơn vị là $400 + \dots$ Số còn thiếu là bao nhiêu?
\shortans{8}
\loigiai{
$408 = 400 + 8$.
}
\end{ex}

\begin{ex}
Số 690 viết thành tổng các trăm và chục là $600 + \dots$ Số còn thiếu là bao nhiêu?
\shortans{90}
\loigiai{
$690 = 600 + 90$.
}
\end{ex}

\begin{ex}
Tính tổng: $200 + 50 + 6$.
\shortans{256}
\loigiai{
Ta có: $200 + 50 + 6 = 256$.
}
\end{ex}

\begin{ex}
Tính tổng: $400 + 80 + 3$.
\shortans{483}
\loigiai{
Ta có: $400 + 80 + 3 = 483$.
}
\end{ex}

\begin{ex}
Tính tổng: $600 + 70 + 9$.
\shortans{679}
\loigiai{
Ta có: $600 + 70 + 9 = 679$.
}
\end{ex}

\begin{ex}
Tính tổng: $800 + 30$.
\shortans{830}
\loigiai{
Ta có: $800 + 30 = 830$.
}
\end{ex}

\begin{ex}
Tính tổng: $900 + 5$.
\shortans{905}
\loigiai{
Ta có: $900 + 5 = 905$.
}
\end{ex}

\begin{ex}
Tính tổng: $700 + 40 + 2$.
\shortans{742}
\loigiai{
Ta có: $700 + 40 + 2 = 742$.
}
\end{ex}

\begin{ex}
Tính tổng: $100 + 90 + 9$.
\shortans{199}
\loigiai{
Ta có: $100 + 90 + 9 = 199$.
}
\end{ex}

\begin{ex}
Tính tổng: $300 + 8$.
\shortans{308}
\loigiai{
Ta có: $300 + 8 = 308$.
}
\end{ex}

\begin{ex}
Tính tổng: $500 + 50$.
\shortans{550}
\loigiai{
Ta có: $500 + 50 = 550$.
}
\end{ex}

\begin{ex}
Một số có tổng các trăm, chục, đơn vị là $600 + 30 + 4$. Số đó là số nào?
\shortans{634}
\loigiai{
Số đó là 634.
}
\end{ex}

\begin{ex}
Một số có tổng các trăm, chục, đơn vị là $800 + 0 + 9$. Số đó là số nào?
\shortans{809}
\loigiai{
Số đó là 809.
}
\end{ex}

\begin{ex}
Số 888 viết thành tổng: $800 + 80 + \dots$ Số cần điền là bao nhiêu?
\shortans{8}
\loigiai{
Số cần điền là 8.
}
\end{ex}

\begin{ex}
Giá trị của chữ số 5 trong số 524 là bao nhiêu?
\shortans{500}
\loigiai{
Chữ số 5 ở hàng trăm có giá trị là 500.
}
\end{ex}

\begin{ex}
Giá trị của chữ số 7 trong số 376 là bao nhiêu?
\shortans{70}
\loigiai{
Chữ số 7 ở hàng chục có giá trị là 70.
}
\end{ex}

\begin{ex}
Giá trị của chữ số 9 trong số 819 là bao nhiêu?
\shortans{9}
\loigiai{
Chữ số 9 ở hàng đơn vị có giá trị là 9.
}
\end{ex}

\begin{ex}
Hiệu giữa giá trị chữ số 6 ở hàng trăm và chữ số 6 ở hàng chục trong số 660 là bao nhiêu?
\shortans{540}
\loigiai{
Giá trị: $600 - 60 = 540$.
}
\end{ex}

\begin{ex}
Một số có chữ số hàng trăm là 4, hàng chục gấp đôi hàng trăm (bằng 8), hàng đơn vị bằng 1. Số đó là số nào?
\shortans{481}
\loigiai{
Số đó là 481.
}
\end{ex}

\begin{ex}
Một số có chữ số hàng trăm là 5, hàng chục là 0, hàng đơn vị lớn nhất có một chữ số (bằng 9). Số đó là số nào?
\shortans{509}
\loigiai{
Số đó là 509.
}
\end{ex}

\begin{ex}
Điền số thích hợp: $754 - 54 = \dots$
\shortans{700}
\loigiai{
Số còn lại là 700.
}
\end{ex}

\begin{ex}
Điền số thích hợp: $836 - 800 = \dots$
\shortans{36}
\loigiai{
Số còn lại là 36.
}
\end{ex}
